\documentclass[11pt,a4paper,UTF8]{ctexart}
\usepackage{geometry}
\geometry{left=18mm,right=18mm,top=24mm,bottom=24mm,headheight=22pt,footskip=12mm}
\usepackage{amsmath,amssymb,mathtools,bm,esint,extarrows}
\usepackage{xcolor}
\usepackage{tcolorbox}
\tcbuselibrary{skins,breakable}
\usepackage{fancyhdr}
\usepackage{titlesec}
\usepackage{hyperref}
\usepackage{tikz}
\usepackage{booktabs}
\usepackage{tabularx}

% --- Color Palette (Purple & DeepPink Academic Theme) ---
\definecolor{DeepPurple}{HTML}{4C1D95}
\definecolor{TitlePurple}{HTML}{581C87}
\definecolor{SubPurple}{HTML}{6B21A8}
\definecolor{DeepPink}{HTML}{FF1493}
\definecolor{LightPinkBg}{HTML}{FFF5F8}
\definecolor{LightPurpleBg}{HTML}{F8F5FF}
\definecolor{GoldAccent}{HTML}{D97706}
\definecolor{DarkText}{HTML}{1F2937}

\hypersetup{
    colorlinks=true,
    linkcolor=TitlePurple,
    citecolor=DeepPink,
    urlcolor=SubPurple,
    pdfborder={0 0 0}
}

% --- Typography & Spacing ---
\linespread{1.3}
\setlength{\parskip}{0.35em plus 0.1em minus 0.1em}
\setlength{\parindent}{0pt}

% --- Section Titles Styling ---
\titleformat{\section}
  {\Large\bfseries\color{TitlePurple}}{\thesection}{1em}{}
  [\vspace{1mm}{\color{TitlePurple!30}\hrule height 1pt}]
\titleformat{\subsection}
  {\large\bfseries\color{SubPurple}}{\thesubsection}{0.8em}{}
\titleformat{\subsubsection}
  {\normalsize\bfseries\color{DeepPurple}}{\thesubsubsection}{0.6em}{}

% --- tcolorbox Environments ---
% 1. Theorem / Formula / Method / Analysis Box (Deep Purple)
\newtcolorbox{thmbox}[1][]{
    enhanced,
    breakable,
    colback=LightPurpleBg,
    colframe=TitlePurple,
    arc=3pt,
    boxrule=0pt,
    leftrule=3.5pt,
    left=10pt,
    right=10pt,
    top=8pt,
    bottom=8pt,
    title=#1,
    coltitle=white,
    colbacktitle=TitlePurple,
    fonttitle=\bfseries\small,
    attach boxed title to top left={yshift=-2mm, xshift=4mm},
    boxed title style={boxrule=0pt, arc=2pt}
}

% 2. Solution / Formula / Derivation Box (DeepPink)
\newtcolorbox{solbox}[1][]{
    enhanced,
    breakable,
    colback=LightPinkBg,
    colframe=DeepPink,
    arc=3pt,
    boxrule=0pt,
    leftrule=3.5pt,
    left=10pt,
    right=10pt,
    top=8pt,
    bottom=8pt,
    title=#1,
    coltitle=white,
    colbacktitle=DeepPink,
    fonttitle=\bfseries\small,
    attach boxed title to top left={yshift=-2mm, xshift=4mm},
    boxed title style={boxrule=0pt, arc=2pt}
}

% 3. Learning Goals / Summary Box
\newtcolorbox{targetbox}[1][]{
    enhanced,
    breakable,
    colback=LightPurpleBg,
    colframe=DeepPurple,
    arc=3pt,
    boxrule=1pt,
    left=10pt,
    right=10pt,
    top=8pt,
    bottom=8pt,
    title=#1,
    coltitle=white,
    colbacktitle=DeepPurple,
    fonttitle=\bfseries\small,
    attach boxed title to top left={yshift=-2mm, xshift=4mm},
    boxed title style={boxrule=0pt, arc=2pt}
}

\begin{document}

% ------------------- 讲义封面 (Cover Page) -------------------
\begin{titlepage}
\thispagestyle{empty}
\begin{tikzpicture}[remember picture,overlay]
    \fill[DeepPurple] (current page.north west) rectangle ([yshift=-7cm]current page.north east);
    \draw[DeepPink, line width=3.5pt] ([yshift=-7cm]current page.north west) -- ([yshift=-7cm]current page.north east);
    \draw[GoldAccent, line width=1pt] ([yshift=-7.12cm]current page.north west) -- ([yshift=-7.12cm]current page.north east);
    
    \node[anchor=north east, opacity=0.12, text=white] at ([xshift=-1.5cm, yshift=-0.5cm]current page.north east) {
        \fontsize{70}{70}\selectfont\bfseries 2027
    };
\end{tikzpicture}

\vspace*{0.8cm}
\begin{center}
    {\color{white}\fontsize{26}{32}\selectfont\bfseries 2027 考研数学(二) 核心讲义}\\[0.6cm]
    {\color{white}\fontsize{13}{18}\selectfont\bfseries 全国硕士研究生招生考试 · 统考数学(二) 核心精讲全书}\\[1.5cm]
\end{center}

\vspace{3.5cm}
\begin{center}
    {\color{GoldAccent}\large\bfseries $\blacktriangleright$ 基础强化 · 核心题解 · 考点深水区 $\blacktriangleleft$}\\[0.8cm]
    {\color{TitlePurple}\fontsize{24}{28}\selectfont\bfseries 第九章 一元积分学的计算}\\[0.6cm]
    {\color{DeepPink}\large\bfseries —— 核心精讲大例题与题解三步法 ——}\\[1.5cm]
\end{center}

\vfill

\begin{center}
\begin{tcolorbox}[
    colback=white,
    colframe=DeepPurple,
    width=0.88\textwidth,
    arc=4pt,
    boxrule=1.2pt,
    halign=center,
    drop shadow=gray!30
]
    \vspace{3mm}
    {\bfseries\large 编著团队：EscoffierZhou}\\[2.5mm]
    {\bfseries\color{TitlePurple} 目标院校：中国科学院大学 (UCAS) 计算机专硕 (22408)}\\[2.5mm]
    {\small\color{gray} 备考铁律：手算真题数字 · 错题白纸重做 · 408 客观题为王}\\[2.5mm]
    {\small 编制版本：2027 届首发版 · \today}
    \vspace{2mm}
\end{tcolorbox}
\end{center}
\vspace{0.8cm}
\end{titlepage}

% ------------------- 目录与页面主体配置 -------------------
\pagestyle{fancy}
\fancyhf{}
\fancyhead[L]{\small\color{gray}2027 考研数学(二) · 核心讲义系列}
\fancyhead[R]{\small\color{TitlePurple}\nouppercase{\leftmark}}
\fancyfoot[C]{\small\color{gray}— 第 \thepage\ 页 —}
\fancyfoot[R]{\small\color{gray}UCAS 22408}
\renewcommand{\headrulewidth}{0.5pt}

\pagenumbering{Roman}
\tableofcontents
\newpage
\pagenumbering{arabic}


\section{9.例题}


\subsection{例题 9.1: 凑微分法与反正弦公式}

求不定积分 $\int \frac{\sqrt{x}}{\sqrt{4-x^3}}\,\mathrm{d}x$。

\textbf{\color{DeepPurple}【思路点拨】} 观察分母根号下为 $4 - x^3 = 2^2 - (x^{3/2})^2$，分子 $\sqrt{x}$ 恰好为内层函数 $x^{3/2}$ 导数的主部，通过调整系数凑微分化为标准反正弦积分


\begin{solbox}[分步推导与答题规范]
\textbf{[1]} 识别内层复合函数与凑微分：

分母根号内可化为平方差结构：


\begin{equation*}
4 - x^3 = 2^2 - \left(x^{\frac{3}{2}}\right)^2
\end{equation*}


计算导数：$\mathrm{d}\left(x^{\frac{3}{2}}\right) = \frac{3}{2}x^{\frac{1}{2}}\,\mathrm{d}x = \frac{3}{2}\sqrt{x}\,\mathrm{d}x$，故分子部分可凑微分为：


\begin{equation*}
\sqrt{x}\,\mathrm{d}x = \frac{2}{3}\,\mathrm{d}\left(x^{\frac{3}{2}}\right)
\end{equation*}


\textbf{[2]} 套用标准积分公式：


\begin{equation*}
\int \frac{\sqrt{x}}{\sqrt{4-x^3}}\,\mathrm{d}x = \frac{2}{3}\int \frac{\mathrm{d}\left(x^{\frac{3}{2}}\right)}{\sqrt{2^2 - \left(x^{\frac{3}{2}}\right)^2}}
\end{equation*}


由基本积分公式 $\int \frac{\mathrm{d}u}{\sqrt{a^2-u^2}} = \arcsin\frac{u}{a} + C$，即得：


\begin{equation*}
= \frac{2}{3}\arcsin\frac{x^{\frac{3}{2}}}{2} + C
\end{equation*}

\end{solbox}


\subsection{例题 9.2: 复杂指数复合与内层导数凑微分}

求不定积分 $\int e^{\frac{\sin\theta}{\cos\theta+\sin\theta}}\frac{1}{(\cos\theta+\sin\theta)^2}\,\mathrm{d}\theta$。

\textbf{\color{DeepPurple}【思路点拨】} 遇到复杂指数复合函数 $e^{u(\theta)}$，首先考察指数内层函数 $u(\theta) = \frac{\sin\theta}{\cos\theta+\sin\theta}$ 的导数，发现被积函数外部因式恰好与之完全吻合


\begin{solbox}[分步推导与答题规范]
\textbf{[1]} 考察指数内层函数导数：

令 $u = \frac{\sin\theta}{\cos\theta+\sin\theta}$，应用商的求导法则求微分：


\begin{equation*}
\mathrm{d}u = \frac{(\sin\theta)'(\cos\theta+\sin\theta) - \sin\theta(\cos\theta+\sin\theta)'}{(\cos\theta+\sin\theta)^2}\,\mathrm{d}\theta
\end{equation*}



\begin{equation*}
= \frac{\cos\theta(\cos\theta+\sin\theta) - \sin\theta(-\sin\theta+\cos\theta)}{(\cos\theta+\sin\theta)^2}\,\mathrm{d}\theta
\end{equation*}



\begin{equation*}
= \frac{\cos^2\theta + \sin\theta\cos\theta + \sin^2\theta - \sin\theta\cos\theta}{(\cos\theta+\sin\theta)^2}\,\mathrm{d}\theta = \frac{1}{(\cos\theta+\sin\theta)^2}\,\mathrm{d}\theta
\end{equation*}


\textbf{[2]} 直接凑微分计算：

原积分化为最简指数积分：


\begin{equation*}
\int e^{\frac{\sin\theta}{\cos\theta+\sin\theta}}\frac{1}{(\cos\theta+\sin\theta)^2}\,\mathrm{d}\theta = \int e^u\,\mathrm{d}u = e^u + C = e^{\frac{\sin\theta}{\cos\theta+\sin\theta}} + C
\end{equation*}

\end{solbox}


\subsection{例题 9.3: 三角代换与直角三角形辅助角回代}

求不定积分 $\int \sqrt{a^2-x^2}\,\mathrm{d}x$ ($a > 0$)。

\textbf{\color{DeepPurple}【思路点拨】} 根式为 $\sqrt{a^2-x^2}$ 型，采用正弦代换 $x = a\sin t$ 消去根号，利用三角恒等式降幂积分，最后通过直角三角形几何比例准确回代


\begin{solbox}[分步推导与答题规范]
\textbf{[1]} 变量代换消根：

令 $x = a\sin t \quad \left(-\frac{\pi}{2} \leqslant t \leqslant \frac{\pi}{2}\right)$，则 $\mathrm{d}x = a\cos t\,\mathrm{d}t$。

此时 $\sqrt{a^2-x^2} = \sqrt{a^2(1-\sin^2 t)} = a|\cos t| = a\cos t$。

原式化为：


\begin{equation*}
\int \sqrt{a^2-x^2}\,\mathrm{d}x = \int a\cos t \cdot a\cos t\,\mathrm{d}t = a^2 \int \cos^2 t\,\mathrm{d}t
\end{equation*}


\textbf{[2]} 倍角公式降幂积分：


\begin{equation*}
a^2 \int \frac{1+\cos 2t}{2}\,\mathrm{d}t = \frac{a^2}{2}\left(t + \frac{1}{2}\sin 2t\right) + C = \frac{a^2}{2}t + \frac{a^2}{2}\sin t\cos t + C
\end{equation*}


\textbf{[3]} 辅助直角三角形回代：

由 $x = a\sin t$，得 $\sin t = \frac{x}{a}$，$t = \arcsin\frac{x}{a}$。

画直角三角形(锐角为 $t$，对边为 $x$，斜边为 $a$，邻边为 $\sqrt{a^2-x^2}$)，即得 $\cos t = \frac{\sqrt{a^2-x^2}}{a}$。

代入整理即得：


\begin{equation*}
= \frac{a^2}{2}\arcsin\frac{x}{a} + \frac{a^2}{2}\cdot \frac{x}{a}\cdot \frac{\sqrt{a^2-x^2}}{a} + C = \frac{x}{2}\sqrt{a^2-x^2} + \frac{a^2}{2}\arcsin\frac{x}{a} + C
\end{equation*}

\end{solbox}


\subsection{例题 9.4: 根式代换与分部积分法结合}

求不定积分 $\int \frac{xe^x}{\sqrt{e^x-1}}\,\mathrm{d}x$。

\textbf{\color{DeepPurple}【思路点拨】} 式中含有复杂根式 $\sqrt{e^x-1}$，直接令 $t = \sqrt{e^x-1}$ 彻底消除根号，反解出 $x$ 与 $\mathrm{d}x$ 后转化为对数函数的分部积分


\begin{solbox}[分步推导与答题规范]
\textbf{[1]} 根式换元消根：

令 $t = \sqrt{e^x-1} \ (t > 0)$，则 $t^2 = e^x - 1 \implies e^x = t^2 + 1$。

两边取对数得 $x = \ln(t^2+1)$，求微分得 $\mathrm{d}x = \frac{2t}{t^2+1}\,\mathrm{d}t$。

代入原式得：


\begin{equation*}
\int \frac{\ln(t^2+1)\cdot(t^2+1)}{t} \cdot \frac{2t}{t^2+1}\,\mathrm{d}t = 2\int \ln(t^2+1)\,\mathrm{d}t
\end{equation*}


\textbf{[2]} 对数函数分部积分：


\begin{equation*}
2\int \ln(t^2+1)\,\mathrm{d}t = 2t\ln(t^2+1) - 2\int t\,\mathrm{d}\ln(t^2+1)
\end{equation*}



\begin{equation*}
= 2t\ln(t^2+1) - 2\int t \cdot \frac{2t}{t^2+1}\,\mathrm{d}t = 2t\ln(t^2+1) - 4\int \frac{t^2+1-1}{t^2+1}\,\mathrm{d}t
\end{equation*}



\begin{equation*}
= 2t\ln(t^2+1) - 4\int \left(1 - \frac{1}{t^2+1}\right)\,\mathrm{d}t = 2t\ln(t^2+1) - 4t + 4\arctan t + C
\end{equation*}


\textbf{[3]} 回代原变量 $t = \sqrt{e^x-1}$：

注意 $t^2+1 = e^x \implies \ln(t^2+1) = x$，代入得：


\begin{equation*}
= 2x\sqrt{e^x-1} - 4\sqrt{e^x-1} + 4\arctan\sqrt{e^x-1} + C
\end{equation*}

\end{solbox}


\subsection{例题 9.5: 三角正切代换与循环型分部积分}

求不定积分 $\int \frac{xe^{\arctan x}}{(1+x^2)^{\frac{3}{2}}}\,\mathrm{d}x$。

\textbf{\color{DeepPurple}【思路点拨】} 式中含有 $\arctan x$ 与 $(1+x^2)$，令 $x = \tan t$ 进行三角代换，被积函数将大幅精简为指数与正弦的经典循环积分


\begin{solbox}[分步推导与答题规范]
\textbf{[1]} 三角换元化简：

令 $x = \tan t \quad \left(-\frac{\pi}{2} < t < \frac{\pi}{2}\right)$，则 $t = \arctan x$，$\mathrm{d}x = \sec^2 t\,\mathrm{d}t$。

分母化为：$(1+x^2)^{\frac{3}{2}} = (1+\tan^2 t)^{\frac{3}{2}} = (\sec^2 t)^{\frac{3}{2}} = \sec^3 t$。

代入原式：


\begin{equation*}
\int \frac{\tan t \cdot e^t}{\sec^3 t} \cdot \sec^2 t\,\mathrm{d}t = \int \frac{\tan t \cdot e^t}{\sec t}\,\mathrm{d}t = \int \frac{\frac{\sin t}{\cos t} \cdot e^t}{\frac{1}{\cos t}}\,\mathrm{d}t = \int e^t \sin t\,\mathrm{d}t
\end{equation*}


\textbf{[2]} 指数三角乘积的循环分部积分：

记 $I = \int e^t \sin t\,\mathrm{d}t$：


\begin{equation*}
I = \int \sin t\,\mathrm{d}(e^t) = e^t \sin t - \int e^t \cos t\,\mathrm{d}t
\end{equation*}



\begin{equation*}
= e^t \sin t - \int \cos t\,\mathrm{d}(e^t) = e^t \sin t - \left(e^t \cos t - \int e^t (-\sin t)\,\mathrm{d}t\right)
\end{equation*}



\begin{equation*}
= e^t (\sin t - \cos t) - I
\end{equation*}


移项解得：$I = \frac{e^t}{2}(\sin t - \cos t) + C$。

\textbf{[3]} 利用辅助三角形回代 $x$：

由 $x = \tan t$，构造直角三角形(对边为 $x$，邻边为 $1$，斜边为 $\sqrt{1+x^2}$)，得：


\begin{equation*}
\sin t = \frac{x}{\sqrt{1+x^2}}, \qquad \cos t = \frac{1}{\sqrt{1+x^2}}
\end{equation*}


代入即得最终原函数：


\begin{equation*}
= \frac{e^{\arctan x}}{2}\left(\frac{x}{\sqrt{1+x^2}} - \frac{1}{\sqrt{1+x^2}}\right) + C = \frac{(x-1)e^{\arctan x}}{2\sqrt{1+x^2}} + C
\end{equation*}

\end{solbox}


\subsection{例题 9.6: 复合函数关系还原与分部积分}

已知 $f(\ln x) = \frac{\ln(1+x)}{x}$，求不定积分 $\int f(x)\,\mathrm{d}x$。

\textbf{\color{DeepPurple}【思路点拨】} 先利用自变量换元还原出 $f(x)$ 的解析式，再利用分部积分法将对数项转化为有理指数分式积分


\begin{solbox}[分步推导与答题规范]
\textbf{[1]} 还原函数解析式：

令 $u = \ln x$，则 $x = e^u$。代入已知关系得：


\begin{equation*}
f(u) = \frac{\ln(1+e^u)}{e^u}
\end{equation*}


因此函数解析式为：$f(x) = \frac{\ln(1+e^x)}{e^x} = e^{-x}\ln(1+e^x)$。

\textbf{[2]} 分部积分计算：


\begin{equation*}
\int f(x)\,\mathrm{d}x = \int \ln(1+e^x)\,\mathrm{d}(-e^{-x}) = -e^{-x}\ln(1+e^x) + \int e^{-x}\,\mathrm{d}\ln(1+e^x)
\end{equation*}



\begin{equation*}
= -e^{-x}\ln(1+e^x) + \int e^{-x} \cdot \frac{e^x}{1+e^x}\,\mathrm{d}x = -e^{-x}\ln(1+e^x) + \int \frac{1}{1+e^x}\,\mathrm{d}x
\end{equation*}


\textbf{[3]} 求解剩余分式积分：

对 $\int \frac{1}{1+e^x}\,\mathrm{d}x$，分子分母同乘以 $e^{-x}$：


\begin{equation*}
\int \frac{e^{-x}}{e^{-x}+1}\,\mathrm{d}x = -\int \frac{\mathrm{d}(e^{-x}+1)}{e^{-x}+1} = -\ln(e^{-x}+1) = -\ln\left(\frac{1+e^x}{e^x}\right) = x - \ln(1+e^x)
\end{equation*}


合并各项得：


\begin{equation*}
= -e^{-x}\ln(1+e^x) + x - \ln(1+e^x) + C = x - (1+e^{-x})\ln(1+e^x) + C
\end{equation*}

\end{solbox}


\subsection{例题 9.7: 完全平方展开与分部积分对消}

求不定积分 $\int e^{2x}(\tan x + 1)^2\,\mathrm{d}x$。

\textbf{\color{DeepPurple}【思路点拨】} 将 $(\tan x + 1)^2$ 展开并利用三角恒等式 $\tan^2 x + 1 = \sec^2 x$，分拆为两项后对正割平方项分部积分，发现后一项恰好被完全对消


\begin{solbox}[分步推导与答题规范]
\textbf{[1]} 三角展开与结构分解：


\begin{equation*}
(\tan x + 1)^2 = \tan^2 x + 2\tan x + 1 = (\tan^2 x + 1) + 2\tan x = \sec^2 x + 2\tan x
\end{equation*}


原积分拆为两部分：


\begin{equation*}
\int e^{2x}(\tan x + 1)^2\,\mathrm{d}x = \int e^{2x}\sec^2 x\,\mathrm{d}x + 2\int e^{2x}\tan x\,\mathrm{d}x
\end{equation*}


\textbf{[2]} 对第一项实施分部积分：

由于 $\sec^2 x\,\mathrm{d}x = \mathrm{d}(\tan x)$：


\begin{equation*}
\int e^{2x}\sec^2 x\,\mathrm{d}x = \int e^{2x}\,\mathrm{d}(\tan x) = e^{2x}\tan x - \int \tan x\,\mathrm{d}(e^{2x}) = e^{2x}\tan x - 2\int e^{2x}\tan x\,\mathrm{d}x
\end{equation*}


\textbf{[3]} 两项相加完全抵消：


\begin{equation*}
\int e^{2x}(\tan x + 1)^2\,\mathrm{d}x = \left(e^{2x}\tan x - 2\int e^{2x}\tan x\,\mathrm{d}x\right) + 2\int e^{2x}\tan x\,\mathrm{d}x = e^{2x}\tan x + C
\end{equation*}

\end{solbox}


\subsection{例题 9.8: 真分式部分分式拆分与待定系数赋值法}

将有理真分式 $\frac{x^2+1}{(x-1)^2(x+2)}$ 分解为部分分式之和。

\textbf{\color{DeepPurple}【思路点拨】} 分母含重根 $(x-1)^2$ 与单根 $(x+2)$，设出标准待定形式后，运用极点覆盖赋值法快速确定最高重根项与单根项系数，再代入特殊值确定剩余未知数


\begin{solbox}[分步推导与答题规范]
\textbf{[1]} 设立待定分式：


\begin{equation*}
\frac{x^2+1}{(x-1)^2(x+2)} = \frac{A}{x-1} + \frac{B}{(x-1)^2} + \frac{C}{x+2}
\end{equation*}


\textbf{[2]} 极点覆盖赋值速算：

- 两端同乘 $(x-1)^2$ 并令 $x=1$：


\begin{equation*}
B = \left. \frac{x^2+1}{x+2} \right|_{x=1} = \frac{1+1}{1+2} = \frac{2}{3}
\end{equation*}


- 两端同乘 $(x+2)$ 并令 $x=-2$：


\begin{equation*}
C = \left. \frac{x^2+1}{(x-1)^2} \right|_{x=-2} = \frac{(-2)^2+1}{(-2-1)^2} = \frac{5}{9}
\end{equation*}


\textbf{[3]} 赋特殊值确定 $A$：

令 $x=0$ 代入恒等式两端：


\begin{equation*}
\frac{0+1}{(-1)^2(2)} = \frac{A}{-1} + \frac{B}{1} + \frac{C}{2} \implies \frac{1}{2} = -A + \frac{2}{3} + \frac{5}{18}
\end{equation*}


解得：$A = \frac{2}{3} + \frac{5}{18} - \frac{1}{2} = \frac{12+5-9}{18} = \frac{8}{18} = \frac{4}{9}$。

\textbf{[4]} 结论：


\begin{equation*}
\frac{x^2+1}{(x-1)^2(x+2)} = \frac{4}{9(x-1)} + \frac{2}{3(x-1)^2} + \frac{5}{9(x+2)}
\end{equation*}

\end{solbox}


\subsection{例题 9.9: 因式分解与有理函数积分}

求不定积分 $\int \frac{x}{x^3-x^2+x-1}\,\mathrm{d}x$。

\textbf{\color{DeepPurple}【思路点拨】} 对分母进行分组因式分解得 $(x-1)(x^2+1)$，拆分为一次项与不可约二次项，通过待定系数求解各部分原函数


\begin{solbox}[分步推导与答题规范]
\textbf{[1]} 分母因式分解与拆项：


\begin{equation*}
x^3 - x^2 + x - 1 = x^2(x-1) + (x-1) = (x-1)(x^2+1)
\end{equation*}


设 $\frac{x}{(x-1)(x^2+1)} = \frac{A}{x-1} + \frac{Bx+C}{x^2+1}$。

通分得：$x = A(x^2+1) + (Bx+C)(x-1)$。

\textbf{[2]} 待定系数求解：

- 令 $x=1$：$1 = A(1+1) \implies A = \frac{1}{2}$。
- 令 $x=0$：$0 = A - C \implies C = A = \frac{1}{2}$。
- 比较 $x^2$ 项系数：$0 = A + B \implies B = -A = -\frac{1}{2}$。

\textbf{[3]} 分别积分求解：


\begin{equation*}
\int \frac{x}{x^3-x^2+x-1}\,\mathrm{d}x = \frac{1}{2}\int \frac{1}{x-1}\,\mathrm{d}x - \frac{1}{2}\int \frac{x-1}{x^2+1}\,\mathrm{d}x
\end{equation*}



\begin{equation*}
= \frac{1}{2}\ln|x-1| - \frac{1}{4}\int \frac{2x}{x^2+1}\,\mathrm{d}x + \frac{1}{2}\int \frac{1}{x^2+1}\,\mathrm{d}x
\end{equation*}



\begin{equation*}
= \frac{1}{2}\ln|x-1| - \frac{1}{4}\ln(x^2+1) + \frac{1}{2}\arctan x + C
\end{equation*}

\end{solbox}


\subsection{例题 9.10: 二次不可约分母的配方与导数拆分}

求不定积分 $\int \frac{2x+3}{x^2-x+1}\,\mathrm{d}x$。

\textbf{\color{DeepPurple}【思路点拨】} 分母判别式 $\Delta = 1 - 4 < 0$ 不可约，先将分子拆出分母的导数 $(2x-1)$ 凑对数项，剩余常数通过分母配方转化为标准反正切函数积分


\begin{solbox}[分步推导与答题规范]
\textbf{[1]} 拆解分子为导数项与常数项：

分母 $x^2-x+1$ 的导数为 $2x-1$。改写分子：


\begin{equation*}
2x + 3 = (2x - 1) + 4
\end{equation*}


原积分拆为两部分：


\begin{equation*}
\int \frac{2x+3}{x^2-x+1}\,\mathrm{d}x = \int \frac{2x-1}{x^2-x+1}\,\mathrm{d}x + 4\int \frac{1}{x^2-x+1}\,\mathrm{d}x
\end{equation*}


\textbf{[2]} 分别积分：

- 第一部分凑微分：


\begin{equation*}
\int \frac{2x-1}{x^2-x+1}\,\mathrm{d}x = \int \frac{\mathrm{d}(x^2-x+1)}{x^2-x+1} = \ln(x^2-x+1) + C_1
\end{equation*}


- 第二部分分母配方：


\begin{equation*}
x^2 - x + 1 = \left(x - \frac{1}{2}\right)^2 + \frac{3}{4} = \left(x - \frac{1}{2}\right)^2 + \left(\frac{\sqrt{3}}{2}\right)^2
\end{equation*}



\begin{equation*}
4\int \frac{1}{\left(x - \frac{1}{2}\right)^2 + \left(\frac{\sqrt{3}}{2}\right)^2}\,\mathrm{d}x = 4 \cdot \frac{1}{\frac{\sqrt{3}}{2}}\arctan\left(\frac{x - \frac{1}{2}}{\frac{\sqrt{3}}{2}}\right) = \frac{8}{\sqrt{3}}\arctan\left(\frac{2x-1}{\sqrt{3}}\right) + C_2
\end{equation*}


\textbf{[3]} 合并结论：


\begin{equation*}
= \ln(x^2-x+1) + \frac{8}{\sqrt{3}}\arctan\left(\frac{2x-1}{\sqrt{3}}\right) + C
\end{equation*}

\end{solbox}


\subsection{例题 9.11: 分段函数定积分与绝对值复合}

设 $f(x) = \begin{cases} x^2, & 0 \leqslant x < 1 \\ \\ 2-x, & 1 \leqslant x \leqslant 2 \end{cases}$，计算 $\int_0^2 f(x)\,\mathrm{d}x$ 以及 $\int_0^2 f(|x-1|)\,\mathrm{d}x$。

\textbf{\color{DeepPurple}【思路点拨】} 利用定积分的可加性在分段点处将区间拆开，复合绝对值根据零点分类讨论去绝对值后再分段积分


\begin{solbox}[分步推导与答题规范]
\textbf{[1]} 计算 $\int_0^2 f(x)\,\mathrm{d}x$：


\begin{equation*}
\int_0^2 f(x)\,\mathrm{d}x = \int_0^1 x^2\,\mathrm{d}x + \int_1^2 (2-x)\,\mathrm{d}x
\end{equation*}



\begin{equation*}
= \left[ \frac{x^3}{3} \right]_0^1 + \left[ 2x - \frac{x^2}{2} \right]_1^2 = \frac{1}{3} + \left[(4-2) - \left(2-\frac{1}{2}\right)\right] = \frac{1}{3} + \frac{1}{2} = \frac{5}{6}
\end{equation*}


\textbf{[2]} 分析 $f(|x-1|)$ 的分段表达式：

当 $x \in [0, 2]$ 时，$|x-1| \in [0, 1]$，始终落在第一段定义域 $[0, 1)$ 内：


\begin{equation*}
f(|x-1|) = (|x-1|)^2 = (x-1)^2
\end{equation*}


\textbf{[3]} 计算定积分：


\begin{equation*}
\int_0^2 f(|x-1|)\,\mathrm{d}x = \int_0^2 (x-1)^2\,\mathrm{d}x = \left[ \frac{(x-1)^3}{3} \right]_0^2 = \frac{1^3}{3} - \frac{(-1)^3}{3} = \frac{1}{3} - \left(-\frac{1}{3}\right) = \frac{2}{3}
\end{equation*}

\end{solbox}


\subsection{例题 9.12: 黎曼和定义求极限与奇函数积分}

计算极限 $\lim_{n \to \infty} \sum_{k=1}^n \frac{k}{n^2+k^2}$，并计算定积分 $\int_{-1}^1 \frac{x^3\arcsin x + \sin x}{1+x^2}\,\mathrm{d}x$。

\textbf{\color{DeepPurple}【思路点拨】} 将黎曼和提取 $\frac{1}{n}$ 构造成定积分的标准形式；对称区间定积分拆项利用“奇函数积分为零”直接消除奇项


\begin{solbox}[分步推导与答题规范]
\textbf{[1]} 构造定积分定义求极限：


\begin{equation*}
\lim_{n \to \infty} \sum_{k=1}^n \frac{k}{n^2+k^2} = \lim_{n \to \infty} \frac{1}{n}\sum_{k=1}^n \frac{\frac{k}{n}}{1+\left(\frac{k}{n}\right)^2} = \int_0^1 \frac{x}{1+x^2}\,\mathrm{d}x
\end{equation*}



\begin{equation*}
= \frac{1}{2}\int_0^1 \frac{\mathrm{d}(1+x^2)}{1+x^2} = \frac{1}{2}\left[\ln(1+x^2)\right]_0^1 = \frac{1}{2}\ln 2
\end{equation*}


\textbf{[2]} 对称区间奇偶性拆项：


\begin{equation*}
\int_{-1}^1 \frac{x^3\arcsin x + \sin x}{1+x^2}\,\mathrm{d}x = \int_{-1}^1 \frac{x^3\arcsin x}{1+x^2}\,\mathrm{d}x + \int_{-1}^1 \frac{\sin x}{1+x^2}\,\mathrm{d}x
\end{equation*}


- 考察第二项：$g(x) = \frac{\sin x}{1+x^2}$ 为奇函数，故在对称区间 $[-1, 1]$ 上积分恒等于 0；
- 考察第一项：$(-x)^3\arcsin(-x) = (-x^3)(-\arcsin x) = x^3\arcsin x$，分子为偶函数，分母为偶函数，故全式为偶函数：


\begin{equation*}
= 2\int_0^1 \frac{x^3\arcsin x}{1+x^2}\,\mathrm{d}x
\end{equation*}

\end{solbox}


\subsection{例题 9.13: 偶倍奇零、三角换元与点火公式}

计算定积分 $\int_{-1}^1 \frac{x^4}{\sqrt{1-x^2}}\,\mathrm{d}x$。

\textbf{\color{DeepPurple}【思路点拨】} 积分区间为 $[-1, 1]$，被积函数为偶函数，先化为两倍半区间积分，再作正弦代换直接动用点火公式秒杀


\begin{solbox}[分步推导与答题规范]
\textbf{[1]} 利用对称性化简：

被积函数 $f(x) = \frac{x^4}{\sqrt{1-x^2}}$ 满足 $f(-x) = f(x)$，为连续偶函数：


\begin{equation*}
\int_{-1}^1 \frac{x^4}{\sqrt{1-x^2}}\,\mathrm{d}x = 2\int_0^1 \frac{x^4}{\sqrt{1-x^2}}\,\mathrm{d}x
\end{equation*}


\textbf{[2]} 三角代换“三换”原则：

令 $x = \sin t$，则 $\mathrm{d}x = \cos t\,\mathrm{d}t$。

换积分限：$x=0 \implies t=0$；$x=1 \implies t=\frac{\pi}{2}$。


\begin{equation*}
2\int_0^1 \frac{x^4}{\sqrt{1-x^2}}\,\mathrm{d}x = 2\int_0^{\frac{\pi}{2}} \frac{\sin^4 t}{\cos t} \cdot \cos t\,\mathrm{d}t = 2\int_0^{\frac{\pi}{2}} \sin^4 t\,\mathrm{d}t
\end{equation*}


\textbf{[3]} 套用华里士点火公式：


\begin{equation*}
2\int_0^{\frac{\pi}{2}} \sin^4 t\,\mathrm{d}t = 2 \cdot \left(\frac{3}{4} \cdot \frac{1}{2} \cdot \frac{\pi}{2}\right) = \frac{3\pi}{8}
\end{equation*}

\end{solbox}


\subsection{例题 9.14: 反三角复合函数的几何代换与分部积分}

计算定积分 $\int_0^1 \arcsin\sqrt{1-x^2}\,\mathrm{d}x$。

\textbf{\color{DeepPurple}【思路点拨】} 由于 $x \in [0, 1]$，令 $x = \cos t$，利用三角恒等式将反三角外壳完全剥除，化为初等三角积分


\begin{solbox}[分步推导与答题规范]
\textbf{[1]} 三角换元脱反三角外壳：

令 $x = \cos t \quad \left(0 \leqslant t \leqslant \frac{\pi}{2}\right)$，则 $\mathrm{d}x = -\sin t\,\mathrm{d}t$。

换限：$x=0 \implies t=\frac{\pi}{2}$；$x=1 \implies t=0$。

根式部分：$\sqrt{1-x^2} = \sqrt{1-\cos^2 t} = \sin t$。

因 $t \in [0, \frac{\pi}{2}]$，$\arcsin(\sin t) = t$。

\textbf{[2]} 积分计算：


\begin{equation*}
\int_0^1 \arcsin\sqrt{1-x^2}\,\mathrm{d}x = \int_{\frac{\pi}{2}}^0 t(-\sin t\,\mathrm{d}t) = \int_0^{\frac{\pi}{2}} t\sin t\,\mathrm{d}t
\end{equation*}


分部积分：


\begin{equation*}
= -\int_0^{\frac{\pi}{2}} t\,\mathrm{d}(\cos t) = \left[ -t\cos t \right]_0^{\frac{\pi}{2}} + \int_0^{\frac{\pi}{2}} \cos t\,\mathrm{d}t = (0 - 0) + [\sin t]_0^{\frac{\pi}{2}} = 1
\end{equation*}

\end{solbox}


\subsection{例题 9.15: 三角代换与点火公式标准结合}

计算定积分 $\int_0^1 x^2\sqrt{1-x^2}\,\mathrm{d}x$。

\textbf{\color{DeepPurple}【思路点拨】} 遇到 $x^m (a^2-x^2)^n$ 的定积分，令 $x = \sin t$ 转化为 $\sin^m t \cos^k t$ 的定积分，利用点火公式或三角恒等式求解


\begin{solbox}[分步推导与答题规范]
\textbf{[1]} 换元：

令 $x = \sin t \ (0 \leqslant t \leqslant \frac{\pi}{2})$，$\mathrm{d}x = \cos t\,\mathrm{d}t$。


\begin{equation*}
\int_0^1 x^2\sqrt{1-x^2}\,\mathrm{d}x = \int_0^{\frac{\pi}{2}} \sin^2 t \cdot \cos t \cdot \cos t\,\mathrm{d}t = \int_0^{\frac{\pi}{2}} \sin^2 t \cos^2 t\,\mathrm{d}t
\end{equation*}


\textbf{[2]} 化为单正弦积分套用点火公式：


\begin{equation*}
\sin^2 t \cos^2 t = \sin^2 t(1 - \sin^2 t) = \sin^2 t - \sin^4 t
\end{equation*}


套用点火公式：


\begin{equation*}
\int_0^{\frac{\pi}{2}} \sin^2 t\,\mathrm{d}t = \frac{1}{2} \cdot \frac{\pi}{2} = \frac{\pi}{4}
\end{equation*}



\begin{equation*}
\int_0^{\frac{\pi}{2}} \sin^4 t\,\mathrm{d}t = \frac{3}{4} \cdot \frac{1}{2} \cdot \frac{\pi}{2} = \frac{3\pi}{16}
\end{equation*}


两项相减即得：


\begin{equation*}
= \frac{\pi}{4} - \frac{3\pi}{16} = \frac{\pi}{16}
\end{equation*}

\end{solbox}


\subsection{例题 9.16: 周期函数定积分起点无关性证明}

设连续函数 $f(x)$ 以 $T$ 为周期，证明：对于任意实数 $a$，恒有 $\int_a^{a+T} f(x)\,\mathrm{d}x = \int_0^T f(x)\,\mathrm{d}x$。

\textbf{\color{DeepPurple}【思路点拨】} 将积分区间利用定积分可加性拆分为三段，对包含移动起点的尾段做周期平移代换，证得首尾两段恰好符号相反互相抵消


\begin{solbox}[分步推导与答题规范]
\textbf{[1]} 拆分积分区间：


\begin{equation*}
\int_a^{a+T} f(x)\,\mathrm{d}x = \int_a^0 f(x)\,\mathrm{d}x + \int_0^T f(x)\,\mathrm{d}x + \int_T^{a+T} f(x)\,\mathrm{d}x
\end{equation*}


\textbf{[2]} 对第三段做周期换元：

对 $\int_T^{a+T} f(x)\,\mathrm{d}x$，令 $x = u + T$，则 $\mathrm{d}x = \mathrm{d}u$。

换限：当 $x = T$ 时 $u = 0$；当 $x = a+T$ 时 $u = a$。


\begin{equation*}
\int_T^{a+T} f(x)\,\mathrm{d}x = \int_0^a f(u+T)\,\mathrm{d}u = \int_0^a f(u)\,\mathrm{d}u
\end{equation*}


\textbf{[3]} 首尾抵消得出结论：

将换元结果代入原式：


\begin{equation*}
\int_a^{a+T} f(x)\,\mathrm{d}x = \int_a^0 f(x)\,\mathrm{d}x + \int_0^T f(x)\,\mathrm{d}x + \int_0^a f(u)\,\mathrm{d}u
\end{equation*}


因为 $\int_a^0 f(x)\,\mathrm{d}x = -\int_0^a f(x)\,\mathrm{d}x$，故首尾两项对消，仅余中间项：


\begin{equation*}
= \int_0^T f(x)\,\mathrm{d}x
\end{equation*}


定理得证。
\end{solbox}


\subsection{例题 9.17: 区间再现公式证明与经典对称积分}

证明区间再现公式 $\int_a^b f(x)\,\mathrm{d}x = \int_a^b f(a+b-x)\,\mathrm{d}x$，并计算定积分 $I = \int_0^{\frac{\pi}{2}} \frac{\sqrt{\sin x}}{\sqrt{\sin x} + \sqrt{\cos x}}\,\mathrm{d}x$。

\textbf{\color{DeepPurple}【思路点拨】} 令 $x = a+b-t$ 证明恒等式；将区间再现应用于被积函数，两式相加分子分母完全约分为1


\begin{solbox}[分步推导与答题规范]
\textbf{[1]} 证明区间再现公式：

在 $\int_a^b f(a+b-x)\,\mathrm{d}x$ 中作变量代换，令 $x = a+b-t$，则 $\mathrm{d}x = -\mathrm{d}t$。

换限：$x=a \implies t=b$；$x=b \implies t=a$。


\begin{equation*}
\int_a^b f(a+b-x)\,\mathrm{d}x = \int_b^a f(t)(-\mathrm{d}t) = \int_a^b f(t)\,\mathrm{d}t = \int_a^b f(x)\,\mathrm{d}x
\end{equation*}


\textbf{[2]} 应用区间再现计算定积分：

这里 $a=0, b=\frac{\pi}{2}$，故 $a+b-x = \frac{\pi}{2} - x$。


\begin{equation*}
I = \int_0^{\frac{\pi}{2}} \frac{\sqrt{\sin\left(\frac{\pi}{2}-x\right)}}{\sqrt{\sin\left(\frac{\pi}{2}-x\right)} + \sqrt{\cos\left(\frac{\pi}{2}-x\right)}}\,\mathrm{d}x = \int_0^{\frac{\pi}{2}} \frac{\sqrt{\cos x}}{\sqrt{\cos x} + \sqrt{\sin x}}\,\mathrm{d}x
\end{equation*}


\textbf{[3]} 两式相加除以二：


\begin{equation*}
2I = \int_0^{\frac{\pi}{2}} \frac{\sqrt{\sin x}}{\sqrt{\sin x} + \sqrt{\cos x}}\,\mathrm{d}x + \int_0^{\frac{\pi}{2}} \frac{\sqrt{\cos x}}{\sqrt{\sin x} + \sqrt{\cos x}}\,\mathrm{d}x
\end{equation*}



\begin{equation*}
= \int_0^{\frac{\pi}{2}} \frac{\sqrt{\sin x} + \sqrt{\cos x}}{\sqrt{\sin x} + \sqrt{\cos x}}\,\mathrm{d}x = \int_0^{\frac{\pi}{2}} 1\,\mathrm{d}x = \frac{\pi}{2}
\end{equation*}


故 $I = \frac{\pi}{4}$。
\end{solbox}


\subsection{例题 9.18: 区间再现消 $x$ 技巧与点火公式}

计算定积分 $I = \int_0^\pi \frac{x\sin^9 x}{2-\cos^2 x}\,\mathrm{d}x$。

\textbf{\color{DeepPurple}【思路点拨】} 被积函数中含有孤立的 $x$ 与 $\sin x$ 复合项，利用区间再现公式 $x = \pi - t$ 将 $x$ 彻底消除，再利用区间对称性折半化为点火公式


\begin{solbox}[分步推导与答题规范]
\textbf{[1]} 区间再现消除因子 $x$：

由区间再现公式 $\int_0^\pi x f(\sin x)\,\mathrm{d}x = \frac{\pi}{2}\int_0^\pi f(\sin x)\,\mathrm{d}x$：


\begin{equation*}
I = \frac{\pi}{2}\int_0^\pi \frac{\sin^9 x}{2-\cos^2 x}\,\mathrm{d}x
\end{equation*}


\textbf{[2]} 利用关于 $x = \frac{\pi}{2}$ 的对称性折半：

被积函数关于 $x = \frac{\pi}{2}$ 对称($\sin(\pi-x) = \sin x$, $\cos^2(\pi-x) = \cos^2 x$)：


\begin{equation*}
I = \frac{\pi}{2} \cdot 2\int_0^{\frac{\pi}{2}} \frac{\sin^9 x}{1 + \sin^2 x}\,\mathrm{d}x = \pi \int_0^{\frac{\pi}{2}} \frac{\sin^9 x}{1+\sin^2 x}\,\mathrm{d}x
\end{equation*}


\textbf{[3]} 长除法化简有理式并结合点火公式求解：

令 $\sin^2 x = t$ 或通过代数多项式恒等变形逐项分解降幂，套用华里士公式求得精确解。
\end{solbox}


\subsection{例题 9.19: 二次变限积分次序交换与降阶恒等式}

设函数 $f(x)$ 连续，证明：$\int_0^x \left(\int_0^u f(t)\,\mathrm{d}t\right)\,\mathrm{d}u = \int_0^x (x-t)f(t)\,\mathrm{d}t$。

\textbf{\color{DeepPurple}【思路点拨】} 外层积分变量为 $u$，将内层变限积分作为整体函数实施分部积分，将二重累次积分降阶为单重积分


\begin{solbox}[分步推导与答题规范]
\textbf{[1]} 分部积分设定：

设 $U = \int_0^u f(t)\,\mathrm{d}t$，$\mathrm{d}V = \mathrm{d}u \implies V = u - x$ (技巧: 选择原函数 $V = u - x$ 可使上下限代入时两端同时为 0！)。

亦可直接取常规 $V = u$：


\begin{equation*}
\int_0^x \left(\int_0^u f(t)\,\mathrm{d}t\right)\,\mathrm{d}u = \left[ u \int_0^u f(t)\,\mathrm{d}t \right]_0^x - \int_0^x u \cdot \frac{\mathrm{d}}{\mathrm{d}u}\left(\int_0^u f(t)\,\mathrm{d}t\right)\,\mathrm{d}u
\end{equation*}


\textbf{[2]} 展开计算：


\begin{equation*}
= x \int_0^x f(t)\,\mathrm{d}t - 0 - \int_0^x u f(u)\,\mathrm{d}u
\end{equation*}


\textbf{[3]} 合并为单重积分：

在最后一项中，积分变量 $u$ 为虚拟变量，可记为 $t$；同时将 $x$ 放入第一个积分号内：


\begin{equation*}
= \int_0^x x f(t)\,\mathrm{d}t - \int_0^x t f(t)\,\mathrm{d}t = \int_0^x (x-t)f(t)\,\mathrm{d}t
\end{equation*}


等式得证。此式即著名的柯西连续积分降阶公式(柯西引理)。
\end{solbox}


\subsection{例题 9.20: 变限积分求导与几何切线方程}

求曲线 $y = \int_0^x (t-1)e^{t^2}\,\mathrm{d}t$ 在拐点处的切线方程。

\textbf{\color{DeepPurple}【思路点拨】} 求曲线拐点需先求一阶导数与二阶导数，令二阶导为 0 锁定拐点坐标，再求拐点处的切线斜率并写出点斜式


\begin{solbox}[分步推导与答题规范]
\textbf{[1]} 求一阶导数与二阶导数：


\begin{equation*}
y' = \frac{\mathrm{d}}{\mathrm{d}x}\int_0^x (t-1)e^{t^2}\,\mathrm{d}t = (x-1)e^{x^2}
\end{equation*}


对 $y'$ 再次求导：


\begin{equation*}
y'' = (x-1)'e^{x^2} + (x-1)(e^{x^2})' = e^{x^2} + (x-1)(2x)e^{x^2} = (2x^2 - 2x + 1)e^{x^2}
\end{equation*}


\textbf{[2]} 分析拐点：

考察方程 $2x^2 - 2x + 1 = 0$，判别式 $\Delta = (-2)^2 - 4(2)(1) = 4 - 8 = -4 < 0$。

由于二次式恒正，故 $y'' > 0$ 恒成立，曲线在整个实数域上凹，无拐点！

\textbf{【注】}若题设点为极值点 $x=1$：

当 $x=1$ 时，$y' = 0$，对应极小值点 $(1, \int_0^1 (t-1)e^{t^2}\,\mathrm{d}t)$，此时切线方程为水平直线 $y = \int_0^1 (t-1)e^{t^2}\,\mathrm{d}t$。
\end{solbox}


\subsection{例题 9.21: 绝对值变限积分的洛必达法则求极限}

计算极限 $\lim_{x \to 0} \frac{\int_0^x t\ln(1+t)\,\mathrm{d}t}{x^3}$。

\textbf{\color{DeepPurple}【思路点拨】} 观察为 $\frac{0}{0}$ 型极限，分子为变上限积分，满足洛必达法则条件，求导后利用等价无穷小快速求出极限


\begin{solbox}[分步推导与答题规范]
\textbf{[1]} 判断极限类型：

当 $x \to 0$ 时，分子 $\int_0^0 = 0$，分母 $0^3 = 0$，为标准的 $\frac{0}{0}$ 型未定式。

\textbf{[2]} 应用洛必达法则：


\begin{equation*}
\lim_{x \to 0} \frac{\int_0^x t\ln(1+t)\,\mathrm{d}t}{x^3} = \lim_{x \to 0} \frac{x\ln(1+x)}{3x^2} = \lim_{x \to 0} \frac{\ln(1+x)}{3x}
\end{equation*}


\textbf{[3]} 等价无穷小替换：

当 $x \to 0$ 时，$\ln(1+x) \sim x$：


\begin{equation*}
= \lim_{x \to 0} \frac{x}{3x} = \frac{1}{3}
\end{equation*}

\end{solbox}


\subsection{例题 9.22: 变限积分函数的极值与驻点判别}

设 $F(x) = \int_0^{x^2} (t-1)(t-4)\,\mathrm{d}t$，求 $F(x)$ 的极值点。

\textbf{\color{DeepPurple}【思路点拨】} 应用复合上下限求导准则求出导函数 $F'(x)$，令其为 0 找出所有驻点，通过检查导数在驻点两侧的正负变号判定极值


\begin{solbox}[分步推导与答题规范]
\textbf{[1]} 求复合上限导数：


\begin{equation*}
F'(x) = (x^2-1)(x^2-4) \cdot (x^2)' = 2x(x-1)(x+1)(x-2)(x+2)
\end{equation*}


\textbf{[2]} 求驻点：

令 $F'(x) = 0$，得到 5 个实根：$x_1 = -2, x_2 = -1, x_3 = 0, x_4 = 1, x_5 = 2$。

\textbf{[3]} 符号穿根分析：

导函数最高次为 $2x^5 > 0$，从最右侧正号开始向左奇穿偶回(均为一次单根，全部穿过变号)：

- 当 $x > 2$ 时，$F'(x) > 0$；
- 在 $x=2$ 处由负变正，为\textbf{极小值点}；
- 在 $x=1$ 处由正变负，为\textbf{极大值点}；
- 在 $x=0$ 处由负变正，为\textbf{极小值点}；
- 在 $x=-1$ 处由正变负，为\textbf{极大值点}；
- 在 $x=-2$ 处由负变正，为\textbf{极小值点}。
\end{solbox}


\subsection{例题 9.23: 原函数奇偶性严格证明}

证明：(1) 若 $f(x)$ 为连续奇函数，则其全体原函数均为偶函数；(2) 若 $f(x)$ 为连续偶函数，则其原函数中有且仅有一个是奇函数。

\textbf{\color{DeepPurple}【思路点拨】} 依据原函数与不定积分的关系 $F(x) = \int_0^x f(t)\,\mathrm{d}t + C$，通过定义法作变量代换严格验证 $F(-x)$ 与 $F(x)$ 的代数关系


\begin{solbox}[分步推导与答题规范]
\textbf{[1]} 证明结论(1)：

设 $f(x)$ 在 $\mathbb{R}$ 上连续且 $f(-x) = -f(x)$。任取 $f(x)$ 的一个原函数 $F(x) = \int_0^x f(t)\,\mathrm{d}t + C$。


\begin{equation*}
F(-x) = \int_0^{-x} f(t)\,\mathrm{d}t + C
\end{equation*}


作变量代换令 $t = -u$，则 $\mathrm{d}t = -\mathrm{d}u$；当 $t=0$ 时 $u=0$；当 $t=-x$ 时 $u=x$：


\begin{equation*}
F(-x) = \int_0^x f(-u)(-\mathrm{d}u) + C = \int_0^x [-f(u)](-\mathrm{d}u) + C = \int_0^x f(u)\,\mathrm{d}u + C = F(x)
\end{equation*}


故对于任意常数 $C$，$F(-x) = F(x)$ 恒成立，即奇函数的全体原函数均为偶函数。

\textbf{[2]} 证明结论(2)：

设 $f(x)$ 为连续偶函数 $f(-x) = f(x)$。设原函数为 $G(x) = \int_0^x f(t)\,\mathrm{d}t + C$。

同理作代换 $t = -u$：


\begin{equation*}
G(-x) = \int_0^{-x} f(t)\,\mathrm{d}t + C = \int_0^x f(-u)(-\mathrm{d}u) + C = -\int_0^x f(u)\,\mathrm{d}u + C
\end{equation*}


若 $G(x)$ 为奇函数，必须满足 $G(-x) = -G(x)$：


\begin{equation*}
-\int_0^x f(u)\,\mathrm{d}u + C = -\left(\int_0^x f(u)\,\mathrm{d}u + C\right) = -\int_0^x f(u)\,\mathrm{d}u - C
\end{equation*}


两边对消得 $C = -C \implies C = 0$。

因此，只有当常数 $C=0$ 时，$G(x)$ 才是奇函数，即有且仅有一个原函数是奇函数。
\end{solbox}


\subsection{例题 9.24: 导函数奇偶性与原函数性质综合判定}

设 $f(x)$ 在 $\mathbb{R}$ 上可导，且 $f'(x)$ 为奇函数，满足 $f(0)=0$。证明：$F(x) = \int_0^x f(t)\,\mathrm{d}t$ 为奇函数。

\textbf{\color{DeepPurple}【思路点拨】} 由 $f'(x)$ 为奇函数推出导函数的所有原函数均为偶函数，结合 $f(0)=0$ 确定 $f(x)$ 的奇偶性，进而再对 $F(x)$ 的奇偶性进行传导判定


\begin{solbox}[分步推导与答题规范]
\textbf{[1]} 判定 $f(x)$ 的奇偶性：

已知 $f'(x)$ 为奇函数，由奇函数的原函数都是偶函数可知：


\begin{equation*}
f(x) \text{ 必为偶函数}
\end{equation*}


\textbf{[2]} 判定 $F(x)$ 的奇偶性：

因为 $f(x)$ 为连续偶函数，由例题[9.23]的结论(2)：

以 0 为积分下限的变上限积分 $F(x) = \int_0^x f(t)\,\mathrm{d}t$ 必为奇函数！

命题得证。
\end{solbox}


\subsection{例题 9.25: 变限积分周期性充要条件证明}

设连续函数 $f(x)$ 以 $T$ 为周期，证明：变上限积分 $F(x) = \int_0^x f(t)\,\mathrm{d}t$ 是以 $T$ 为周期的周期函数的充要条件是 $\int_0^T f(t)\,\mathrm{d}t = 0$。

\textbf{\color{DeepPurple}【思路点拨】} 考察周期增量差式 $F(x+T) - F(x)$，利用周期函数积分起点无关性化为一个周期上的定积分，直接建立等价关系


\begin{solbox}[分步推导与答题规范]
\textbf{[1]} 必要性证明($\implies$)：

若 $F(x)$ 以 $T$ 为周期，则对任意 $x$ 都有 $F(x+T) = F(x)$。

特别地，取 $x=0$ 代入：


\begin{equation*}
F(T) = F(0) = \int_0^0 f(t)\,\mathrm{d}t = 0
\end{equation*}


而由定义 $F(T) = \int_0^T f(t)\,\mathrm{d}t$，故必有 $\int_0^T f(t)\,\mathrm{d}t = 0$。

\textbf{[2]} 充分性证明($\impliedby$)：

若 $\int_0^T f(t)\,\mathrm{d}t = 0$。计算周期差式：


\begin{equation*}
F(x+T) - F(x) = \int_0^{x+T} f(t)\,\mathrm{d}t - \int_0^x f(t)\,\mathrm{d}t = \int_x^{x+T} f(t)\,\mathrm{d}t
\end{equation*}


由周期函数定积分起点无关性(例题[9.16])：


\begin{equation*}
\int_x^{x+T} f(t)\,\mathrm{d}t = \int_0^T f(t)\,\mathrm{d}t = 0
\end{equation*}


因此 $F(x+T) - F(x) = 0 \implies F(x+T) = F(x)$。

说明 $F(x)$ 必是以 $T$ 为周期的周期函数。
\end{solbox}


\subsection{例题 9.26: 内部瑕点识别与反常积分拆分}

计算反常积分 $\int_{\frac{1}{2}}^{\frac{3}{2}} \frac{\mathrm{d}x}{\sqrt{|x-x^2|}}$。

\textbf{\color{DeepPurple}【思路点拨】} 绝对值函数在 $x=1$ 处正负变号，且分母在 $x=1$ 处等于 0，说明 $x=1$ 是积分开区间内部的无界瑕点！必须强制拆分成两个瑕积分分别计算


\begin{solbox}[分步推导与答题规范]
\textbf{[1]} 内部瑕点定位与绝对值脱模：

考察被积函数分母 $x - x^2 = x(1-x)$：

- 当 $x \in [\frac{1}{2}, 1)$ 时，$x(1-x) > 0$，故 $|x-x^2| = x-x^2$；
- 当 $x \in (1, \frac{3}{2}]$ 时，$x(1-x) < 0$，故 $|x-x^2| = x^2-x$。

在 $x=1$ 处分母为 0，故 $x=1$ 为内部瑕点，必须强制拆开：


\begin{equation*}
\int_{\frac{1}{2}}^{\frac{3}{2}} \frac{\mathrm{d}x}{\sqrt{|x-x^2|}} = \int_{\frac{1}{2}}^1 \frac{\mathrm{d}x}{\sqrt{x-x^2}} + \int_1^{\frac{3}{2}} \frac{\mathrm{d}x}{\sqrt{x^2-x}}
\end{equation*}


\textbf{[2]} 分别配方计算：

- 第一部分：$x - x^2 = \frac{1}{4} - (x - \frac{1}{2})^2$：


\begin{equation*}
\int_{\frac{1}{2}}^1 \frac{\mathrm{d}x}{\sqrt{\left(\frac{1}{2}\right)^2 - \left(x-\frac{1}{2}\right)^2}} = \left[ \arcsin\left(\frac{x-\frac{1}{2}}{\frac{1}{2}}\right) \right]_{\frac{1}{2}}^1 = \arcsin(1) - \arcsin(0) = \frac{\pi}{2}
\end{equation*}


- 第二部分：$x^2 - x = (x - \frac{1}{2})^2 - \frac{1}{4}$：


\begin{equation*}
\int_1^{\frac{3}{2}} \frac{\mathrm{d}x}{\sqrt{\left(x-\frac{1}{2}\right)^2 - \left(\frac{1}{2}\right)^2}} = \left[ \ln\left| \left(x-\frac{1}{2}\right) + \sqrt{x^2-x} \right| \right]_1^{\frac{3}{2}}
\end{equation*}



\begin{equation*}
= \ln\left( 1 + \frac{\sqrt{3}}{2} \right) - \ln\left( \frac{1}{2} + 0 \right) = \ln(2 + \sqrt{3})
\end{equation*}


\textbf{[3]} 合并结论：


\begin{equation*}
= \frac{\pi}{2} + \ln(2+\sqrt{3})
\end{equation*}

\end{solbox}


\subsection{例题 9.27: 变量平移与倒代换求解反常积分}

计算反常积分 $\int_3^{+\infty} \frac{\mathrm{d}x}{(x-1)^4\sqrt{x^2-2x}}$。

\textbf{\color{DeepPurple}【思路点拨】} 根号下 $x^2-2x = (x-1)^2 - 1$，提示先平移变量令 $u = x-1$，再利用倒代换 $u = \frac{1}{t}$ 消除高次与根式


\begin{solbox}[分步推导与答题规范]
\textbf{[1]} 平移变量：

令 $u = x-1$，则 $\mathrm{d}x = \mathrm{d}u$，根式化为 $\sqrt{u^2-1}$。

换限：$x=3 \implies u=2$；$x \to +\infty \implies u \to +\infty$。

原式化为：


\begin{equation*}
I = \int_2^{+\infty} \frac{\mathrm{d}u}{u^4\sqrt{u^2-1}}
\end{equation*}


\textbf{[2]} 倒代换 $u = \frac{1}{t}$：

令 $u = \frac{1}{t}$，则 $\mathrm{d}u = -\frac{1}{t^2}\,\mathrm{d}t$。

换限：$u=2 \implies t=\frac{1}{2}$；$u \to +\infty \implies t \to 0^+$。


\begin{equation*}
I = \int_{\frac{1}{2}}^0 \frac{-\frac{1}{t^2}\,\mathrm{d}t}{\frac{1}{t^4}\sqrt{\frac{1}{t^2}-1}} = \int_0^{\frac{1}{2}} \frac{t^4 \cdot t}{t^2 \sqrt{1-t^2}}\,\mathrm{d}t = \int_0^{\frac{1}{2}} \frac{t^3}{\sqrt{1-t^2}}\,\mathrm{d}t
\end{equation*}


\textbf{[3]} 凑微分计算：


\begin{equation*}
I = \int_0^{\frac{1}{2}} \frac{t^2}{\sqrt{1-t^2}} \cdot t\,\mathrm{d}t = -\frac{1}{2}\int_0^{\frac{1}{2}} \frac{t^2}{\sqrt{1-t^2}}\,\mathrm{d}(1-t^2)
\end{equation*}


令 $w = \sqrt{1-t^2}$，则 $w^2 = 1-t^2 \implies t^2 = 1-w^2$，$w\,\mathrm{d}w = -t\,\mathrm{d}t$。

换限：$t=0 \implies w=1$；$t=\frac{1}{2} \implies w=\frac{\sqrt{3}}{2}$。


\begin{equation*}
I = \int_{\frac{\sqrt{3}}{2}}^1 (1-w^2)\,\mathrm{d}w = \left[ w - \frac{w^3}{3} \right]_{\frac{\sqrt{3}}{2}}^1 = \left(1 - \frac{1}{3}\right) - \left(\frac{\sqrt{3}}{2} - \frac{3\sqrt{3}}{24}\right)
\end{equation*}



\begin{equation*}
= \frac{2}{3} - \left(\frac{\sqrt{3}}{2} - \frac{\sqrt{3}}{8}\right) = \frac{2}{3} - \frac{3\sqrt{3}}{8}
\end{equation*}

\end{solbox}


\subsection{例题 9.28: 分部积分建立 $\Gamma$ 函数递推关系}

建立递推关系证明：对任意非负整数 $n$，恒有 $I_n = \int_0^{+\infty} x^n e^{-x}\,\mathrm{d}x = n!$。

\textbf{\color{DeepPurple}【思路点拨】} 对被积函数采用分部积分法，求导降低多项式幂次，建立一阶差分递推关系式


\begin{solbox}[分步推导与答题规范]
\textbf{[1]} 分部积分求递推式：


\begin{equation*}
I_n = \int_0^{+\infty} x^n\,\mathrm{d}(-e^{-x}) = \left[ -x^n e^{-x} \right]_0^{+\infty} + \int_0^{+\infty} e^{-x}\,\mathrm{d}(x^n)
\end{equation*}


由极限洛必达法则：对任意自然数 $n$，$\lim_{x \to +\infty} x^n e^{-x} = 0$，且当 $x=0$ 时 $0^n = 0$。

因此端点差值严格为 0：


\begin{equation*}
I_n = 0 + n\int_0^{+\infty} x^{n-1}e^{-x}\,\mathrm{d}x = n I_{n-1} \quad (n \geqslant 1)
\end{equation*}


\textbf{[2]} 计算初值并归纳：


\begin{equation*}
I_0 = \int_0^{+\infty} e^{-x}\,\mathrm{d}x = \left[ -e^{-x} \right]_0^{+\infty} = 0 - (-1) = 1 = 0!
\end{equation*}


由递推公式连续展开：


\begin{equation*}
I_n = n \cdot (n-1) \cdot (n-2) \cdots 1 \cdot I_0 = n!
\end{equation*}


命题得证。
\end{solbox}


\subsection{例题 9.29: $\Gamma$ 函数半整数与高斯分布结合}

设 $f(x) = \begin{cases} \frac{4x^2}{a^3\sqrt{\pi}} e^{-\frac{x^2}{a^2}}, & x > 0 \\ \\ 0, & x \leqslant 0 \end{cases}$ ($a$ 为正常数)，计算定积分 $\int_0^{+\infty} x^2 f(x)\,\mathrm{d}x$。

\textbf{\color{DeepPurple}【思路点拨】} 将 $f(x)$ 代入积分，令 $u = \frac{x}{a}$ 作换元，化为 $\Gamma$ 函数的积分模型，利用 $\Gamma(\frac{5}{2})$ 的半整数递推求得最终结果


\begin{solbox}[分步推导与答题规范]
\textbf{[1]} 列出积分表达式并作变量代换：


\begin{equation*}
\int_0^{+\infty} x^2 f(x)\,\mathrm{d}x = \int_0^{+\infty} x^2 \cdot \frac{4x^2}{a^3\sqrt{\pi}} e^{-\frac{x^2}{a^2}}\,\mathrm{d}x = \frac{4}{a^3\sqrt{\pi}}\int_0^{+\infty} x^4 e^{-\left(\frac{x}{a}\right)^2}\,\mathrm{d}x
\end{equation*}


令 $u = \frac{x}{a}$，则 $x = au$，$\mathrm{d}x = a\,\mathrm{d}u$：


\begin{equation*}
= \frac{4}{a^3\sqrt{\pi}}\int_0^{+\infty} (au)^4 e^{-u^2} (a\,\mathrm{d}u) = \frac{4a^2}{\sqrt{\pi}}\int_0^{+\infty} u^4 e^{-u^2}\,\mathrm{d}u
\end{equation*}


\textbf{[2]} 转化为 $\Gamma$ 函数模型：

改写指数为 $u^{2 \cdot \frac{5}{2} - 1}$：


\begin{equation*}
\int_0^{+\infty} u^4 e^{-u^2}\,\mathrm{d}u = \frac{1}{2} \cdot 2\int_0^{+\infty} u^{2 \cdot \frac{5}{2} - 1} e^{-u^2}\,\mathrm{d}u = \frac{1}{2}\Gamma\left(\frac{5}{2}\right)
\end{equation*}


\textbf{[3]} 代入半整数递推值计算：


\begin{equation*}
\Gamma\left(\frac{5}{2}\right) = \frac{3}{2}\Gamma\left(\frac{3}{2}\right) = \frac{3}{2} \cdot \frac{1}{2}\Gamma\left(\frac{1}{2}\right) = \frac{3}{4}\sqrt{\pi}
\end{equation*}


代入原式：


\begin{equation*}
= \frac{4a^2}{\sqrt{\pi}} \cdot \left(\frac{1}{2} \cdot \frac{3}{4}\sqrt{\pi}\right) = \frac{3}{2}a^2
\end{equation*}

\end{solbox}

\end{document}
